\documentclass[11pt]{article}
\usepackage[margin=25mm]{geometry}
\usepackage{fontspec}
\setmainfont{Latin Modern Roman}
\usepackage{amsmath,amssymb,mathtools}
\usepackage{unicode-math}
\setmathfont{Latin Modern Math}
\usepackage{xurl,hyperref}
\hypersetup{colorlinks=true,urlcolor=blue}
\setcounter{secnumdepth}{0}
\setlength{\parindent}{0pt}
\setlength{\parskip}{5pt}
\setlength{\emergencystretch}{3em}
\newcommand{\sref}[2]{\hyperlink{source:#1:#2}{[S#2]}}
\newcommand{\cataloguescope}[1]{\paragraph{Recorded target.}#1}
\newcommand{\R}{\mathbb R}
\newcommand{\C}{\mathbb C}
\newcommand{\N}{\mathbb N}
\newcommand{\Z}{\mathbb Z}
\newcommand{\Q}{\mathbb Q}
\newcommand{\F}{\mathbb F}
\newcommand{\eps}{\varepsilon}
\begin{document}
\section{Massless and Point-Charge Limits of the Pauli–Fierz $g$-Factor}\label{statement}

\subsection{Definitions and mathematical statement}

\paragraph{Shared notation.}$\mathbb N=\{1,2,\ldots\}$; the standard number fields have their usual meanings. All additional conventions are stated locally.


\paragraph{Definitions and conventions.} The translation-invariant Pauli--Fierz Hamiltonian decomposes into total-momentum fibers $H(p)$ acting on electron-spin and photon Fock space. $m_{\mathrm{ph}}$ regularizes the photon infrared dispersion and $\widehat\varphi(k)$ is the charge form factor providing the ultraviolet cutoff. The $g$-factor measures the spin response relative to the effective orbital magnetic response and is defined by the report's equation (14). The point-charge limit sends $\varphi$ to a delta distribution, equivalently $\widehat\varphi(k)\to(2\pi)^{-3/2}$ in its Fourier convention. The fine-structure constant is $\alpha=e^2/(4\pi)$ in the report's units.

\paragraph{Source scope and qualifications.} The perturbative first correction is not a proof that the full cutoff limits exist. The two cutoff removals, possible order of limits, all-alpha dependence, and the second-order coefficient are separate questions. The g-factor definition is the source's equation (14), including effective-mass normalization.

\paragraph{Catalogue mathematical target.}
\paragraph{Catalogue wording retained for traceability.}
For the Pauli–Fierz Hamiltonian $H(p)$ with photon dispersion $\omega(k)=\sqrt{|k|^2+m_{\mathrm{ph}}^2}$, does its $g$-factor have a finite limit as $m_{\mathrm{ph}}\to0$ and $\widehat\varphi(k)\to(2\pi)^{-3/2}$? If the point-charge limit exists, determine its dependence on the fine-structure constant $\alpha$, including the coefficient of $O(\alpha^2)$ in $$g_\infty=2\left(1+\frac{4\alpha}{3\pi}+O(\alpha^2)\right)$$.

\paragraph{Saved literature context.} The supplied catalogue's assessment dated 2026-08-21 records: Perturbative coefficients and several ultraviolet or point-limit results are known, but the report's full infrared limit, alpha-dependence, and O(alpha squared) questions for the interacting Hamiltonian remain unresolved. The assessment is retained as dated provenance; it is not a new comprehensive literature or open-status certification.
\subsection{Short English statement}
Determine whether the Pauli--Fierz g-factor survives removal of both photon-mass and charge-size cutoffs, and compute its dependence on coupling.
\subsection{Research attempt}
\paragraph{Model, date and investigated scope.}
OpenAI GPT-6 Astra (configured model \texttt{gpt-6-astra}), ultra reasoning; generated 11 October 2026. Provides explicit scalar countertests showing that the known first-order term, even with bounded second-order coefficients, cannot justify a joint cutoff limit.

\paragraph{Sources and literature check.}
The complete supplied definition, including source qualifications and the dated literature assessment, was reviewed. The calculation below is an elementary restricted attempt motivated by [S2]; it does not claim novelty or a fresh literature-status determination.

\paragraph{Context and investigated scope.} The complete supplied definition, including source qualifications and the dated literature assessment, was reviewed. The calculation below is an elementary restricted attempt motivated by [S2]; it does not claim novelty or a fresh literature-status determination.

\paragraph{A cutoff-uniformity stress test.} Let $a=4/(3\pi)$ and introduce positive cutoffs $m,\Lambda$. Consider the illustrative scalar family
\[
 G(\alpha;m,\Lambda)=2\left(1+a\alpha+\alpha^2\frac{m\Lambda}{1+m\Lambda}\right).
\]
For each fixed cutoff pair it has the source's displayed first-order coefficient, and its second-order coefficient is uniformly bounded between zero and two. Nevertheless
\[
 \lim_{m\downarrow0}\lim_{\Lambda\to\infty}G=2(1+a\alpha+\alpha^2),\quad
 \lim_{\Lambda\to\infty}\lim_{m\downarrow0}G=2(1+a\alpha).
\]
Along $m\Lambda=c>0$, the coefficient is $2c/(1+c)$. Thus even bounded second-order remainders and agreement at first order do not imply a joint cutoff limit.

\paragraph{A separate ultraviolet failure mode.} Replacing the bounded fraction by $\log(1+\Lambda)$ preserves the same first-order expansion at every finite cutoff but destroys the ultraviolet limit for each fixed nonzero $\alpha$. This is not a computation of any Pauli--Fierz coefficient: both functions are explicit tests of a proposed inference from the reported perturbative term.

\paragraph{What an actual proof must add.} The source $g$-factor includes effective-mass normalization, so numerator and denominator responses may each depend on both cutoffs. One must control their ratio uniformly and establish that the mixed infrared/ultraviolet limits commute or specify their order. Computing the true second-order coefficient also requires those operator response terms; no value is inferred from the toy family. All coupling dependence, both cutoff removals and the $O(\alpha^2)$ question remain unresolved.

\paragraph{Verification and adversarial audit.}
Checked both iterated limits, fixed-product paths, zero coupling and divergent ultraviolet variant; no toy coefficient was identified with the physical one. Model self-audit only; no independent human proof checking. Prompts A and B were applied to this canonical record and its qualified source scope. Checked both iterated limits, fixed-product paths, zero coupling and divergent ultraviolet variant; no toy coefficient was identified with the physical one. Retained unresolved status and zero original-target progress; Cutoff-uniform control of the actual normalized operator responses and their second-order coefficient is missing.

Coordinator checked both ordered cutoff limits and the path-dependent second-order term of the explicit scalar diagnostic family. This additional check was model-assisted, not independent human verification.

\paragraph{Tools, process and contributions.}
Python 3 was used for file preparation and structural checks; no numerical experiment or theorem prover was used. Contribution prompt A was followed by prompt B. GPT-6 Astra with ultra reasoning drafted the research and performed its stated self-audit; any additional coordinator contribution is identified in the verification and audit. A GPT-6-based Codex coordinator assembled the manuscript and checked repository integration; its exact variant and reasoning configuration were not exposed. No human mathematical contribution, independent human proof check or formal proof certification is supplied. Complete derivations and the consulted source references are included here; private reasoning transcripts are not published.

\paragraph{Final status and first remaining gap.}
\textbf{UNRESOLVED}. Cutoff-uniform control of the actual normalized operator responses and their second-order coefficient is missing.

\section{Completion Estimate}
\noindent\textbf{COMPLETION: 0\%}
The calculation is an auxiliary or restricted attempt and establishes no newly verified portion of the full source target. Cutoff-uniform control of the actual normalized operator responses and their second-order coefficient is missing.

\subsection{Sources}
\begin{itemize}
\item[\hypertarget{source:statement:1}{S1}] ULAM.AI and the UnsolvedMath Contributors, UnsolvedMath, catalogue entry. Catalogue curation under CC BY 4.0; definitions and English exposition added locally.
\url{https://huggingface.co/datasets/ulamai/UnsolvedMath}.
\item[\hypertarget{source:statement:2}{S2}] Mini-Workshop: Attraction to Solitary Waves and Related Aspects of Physics (2008), Oberwolfach Reports, DOI 10.4171/owr/2008/08, pp. 367--400; official report source and mathematical context.
\url{https://doi.org/10.4171/owr/2008/08}.
\item[\hypertarget{source:statement:3}{S3}] Official EMS Press PDF of the same Oberwolfach report; PDF page 29, printed page 395, equations (14)--(16).
\url{https://ems.press/content/serial-article-files/46151}.
\item[\hypertarget{source:statement:4}{S4}] Pauli--Fierz point-limit dynamics; reference in the saved dated catalogue literature review.
\url{https://arxiv.org/abs/math-ph/0407002}.
\item[\hypertarget{source:statement:5}{S5}] ultraviolet-limit analysis; reference in the saved dated catalogue literature review.
\url{https://arxiv.org/abs/2004.06494}.
\end{itemize}
\label{end:statement}
\end{document}
